\chapter{Investigation Numérique (Forensics)}
\label{chap:forensics}

Une fois l'incident détecté et contenu, il faut comprendre \textbf{comment} l'attaquant est entré et \textbf{ce qu'il a fait}. C'est le rôle de l'investigation numérique, ou \textit{Computer Forensics}.

\section{Présentation : Les Objectifs de l'Investigation}

\subsection{Définition}

\begin{boxdef}
\textbf{Investigation Numérique (Forensics) :} Science qui consiste à identifier, préserver, analyser et présenter des preuves numériques de manière à ce qu'elles soient recevables juridiquement.
\end{boxdef}

L'objectif est double :
\begin{itemize}
    \item \textbf{Technique :} Comprendre l'attaque pour reboucher la faille (Root Cause Analysis).
    \item \textbf{Juridique :} Constituer un dossier de preuves pour une éventuelle action en justice ou une assurance.
\end{itemize}

\subsection{Le Principe Fondamental : Ne Rien Altérer}

La règle d'or en forensics est l'intégrité. Toute manipulation de la preuve doit pouvoir être justifiée. Si l'intégrité d'une preuve est douteuse, elle sera rejetée par un tribunal.

% ------------------------------------------------------------------------------
\section{Concepts Clés et Volatilité}
% ------------------------------------------------------------------------------

\subsection{L'Ordre de Volatilité (Order of Volatility)}

Les données numériques ne sont pas toutes égales face au temps. Certaines disparaissent dès que l'on éteint l'ordinateur. Il faut collecter les plus fragiles en premier.

L'ordre standard (RFC 3227) est :
\begin{enumerate}
    \item \textbf{Registres, Cache CPU} (Très volatile).
    \item \textbf{Mémoire Vive (RAM)} : Contient les processus en cours, les mots de passe en clair, les connexions réseaux ouvertes. \textbf{Disparaît à l'extinction.}
    \item \textbf{État du réseau} (Connexions actives, table ARP).
    \item \textbf{Processus en cours}.
    \item \textbf{Disque Dur} : Fichiers, Logs. (Persistant).
    \item \textbf{Configuration distante} (Logs routeurs, pare-feu).
    \item \textbf{Sauvegardes physiques} (Bandes, Archives).
\end{enumerate}

\begin{boxlizbiz}
\textbf{Erreur classique LiZbiZ :}
Un technicien, voyant un poste infecté, l'éteint immédiatement ("Pour stopper le virus").
\textbf{Conséquence Forensique :} La mémoire vive (RAM) est perdue. On ne saura jamais quels programmes le virus lançait ni s'il avait ouvert une connexion distante vers un serveur de commande (C\&C).
\textbf{Bonne pratique :} Faire une "image mémoire" (dump) \textbf{avant} d'éteindre.
\end{boxlizbiz}

\subsection{Chaîne de Custodie (Chain of Custody)}

\begin{boxdef}
\textbf{Chaîne de Custodie :} Document qui retrace l'historique complet d'une preuve depuis sa collecte jusqu'à sa présentation. Qui l'a manipulée, quand, pourquoi, et où elle est stockée.
\end{boxdef}

C'est le document qui garantit l'intégrité de la preuve.

% ------------------------------------------------------------------------------
\section{Jargon de l'Investigateur}
% ------------------------------------------------------------------------------

\subsection{Image Disque et Acquisition}

\begin{itemize}
    \item \textbf{Image Disque (Disk Image) :} Copie bit-à-bit exacte du support de stockage. On travaille sur la copie, jamais sur l'original.
    \item \textbf{Write Blocker :} Matériel physique qui empêche toute écriture sur le disque dur lors de la connexion pour l'acquisition (préservation).
\end{itemize}

\subsection{Hash (Empreinte Cryptographique)}

\begin{boxdef}
\textbf{Hash :} Chaîne de caractères calculée mathématiquement (ex: MD5, SHA-256) qui sert d'empreinte digitale à un fichier. Si un seul bit change, le hash change complètement.
\end{boxdef}

C'est l'outil de preuve d'intégrité.

\begin{boxlizbiz}
\textbf{Procédure LiZbiZ :}
L'investigateur fait une image du disque dur du serveur compromis.
1. Calcul du Hash de l'original : \texttt{a3f2e...}
2. Calcul du Hash de la copie : \texttt{a3f2e...}
Si les hashs sont identiques, la copie est parfaite et recevable en justice.
\end{boxlizbiz}

\subsection{Artéfacts}

\begin{boxdef}
\textbf{Artéfacts :} Traces laissées par l'activité d'un utilisateur ou d'un programme. Ce sont les indices de l'enquête.
\end{boxdef}

Exemples : Historique navigateur, fichiers récemment ouverts (Prefetch), logs Windows (Event Viewer), corbeille, cache DNS.

% ------------------------------------------------------------------------------
\section{Cas LiZbiZ : Réaction Forensique sur le Lab}
% ------------------------------------------------------------------------------

\subsection{Scénario : Analyse Post-Mortem}

\begin{boxlizbiz}
\textbf{Situation :} L'alerte SSH du chapitre précédent a été confirmée. Un attaquant a accédé le serveur ERP via le VLAN IoT. Le serveur est isolé du réseau mais toujours allumé.
\end{boxlizbiz}

\subsection{Travaux Pratiques : Collecte de Preuves}

Les étudiants doivent agir comme des enquêteurs numériques sur la machine virtuelle (Linux) dans GNS3.

\textbf{Étape 1 : Capture de la mémoire (Simulée)}
Sur un système Linux, on peut utiliser des outils comme \textit{LiME} ou \textit{dd} sur \texttt{/dev/mem}. Ici, nous allons simuler la capture de l'état du système.
\begin{itemize}
    \item \textbf{Commande :} Lister les connexions actives suspectes.
    \texttt{ss -tupan} ou \texttt{netstat -tupan}.
    \item \textbf{Objectif :} Voir si l'attaquant est encore connecté ou a laissé une "backdoor" (porte dérobée).
\end{itemize}

\textbf{Étape 2 : Analyse des Logs (Artéfacts)}
L'attaquant s'est connecté en SSH. Où chercher ?
\begin{itemize}
    \item \textbf{Fichier :} \texttt{/var/log/auth.log} (Debian/Ubuntu).
    \item \textbf{Commande :} \texttt{grep "Failed password" /var/log/auth.log} ou \texttt{grep "Accepted password"}.
    \item \textbf{Preuve recherchée :} L'adresse IP de l'attaquant, l'heure de connexion, le compte utilisé.
\end{itemize}

\textbf{Étape 3 : Recherche de fichiers modifiés}
Quels fichiers l'attaquant a-t-il téléchargés ou modifiés ?
\begin{itemize}
    \item \textbf{Commande :} \texttt{find / -mtime -1 -ls}.
    \item \textbf{Explication :} Liste tous les fichiers modifiés dans les dernières 24 heures.
\end{itemize}

\textbf{Étape 4 : Calcul de Hash}
Pour valider l'intégrité d'un fichier suspect trouvé (ex: un script malveillant \texttt{cleanup.sh}).
\begin{itemize}
    \item \textbf{Commande :} \texttt{sha256sum cleanup.sh}.
    \item \textbf{Livrable :} Notez le hash dans le rapport d'investigation.
\end{itemize}

% ==============================================================================
% Fichier : 04_module_incident.tex (Extrait Chapitre 4)
% Description : Continuité d'Activité et Résilience (ISO 22301)
% ==============================================================================